\section{Devin Fusion、Coding Agent Indexで首位を主張}
Cognitionは、Devin FusionがArtificial AnalysisのCoding Agent Indexで68点を獲得したと発表した。Opus 5.5を主担当、GPT-6 Lunaを補助役とする構成で、同じOpus 5.5を用いるClaude Codeより高性能かつタスク当たりの費用が31\%低いという。公式投稿には比較グラフとブログへの案内が含まれ、コミュニティでは費用対効果に注目が集まった。一方、グラフの解釈への疑問もあり、数値の再現性や評価条件は本号で独立確認していない。
\dailyxsource{Cognition (@cognition)}{https://x.com/cognition/status/2112290148205436979}

\section{拡散型LLM「Mercury 3」を発表}
Inceptionは拡散方式のLLM「Mercury 3」を公開した。同社によれば、GPT-6 LunaやClaude Haiku 5.5を複数のベンチマークで上回り、生成速度は毎秒1,000トークン超、コンテキスト長は100万トークンに達する。Mercury 1と比べてモデル規模は40倍超、知能面は3倍という説明もある。APIのほかBasetenやOpenRouterで提供するとされ、高速生成と品質の両立が注目点となった。比較性能や実利用時の遅延・コストは発表側の主張である。
\dailyxsource{Inception (@\_inception\_ai)}{https://x.com/_inception_ai/status/2112236719068369299}

\section{Sunra MCP、画像・動画・音声モデルを横断利用}
Sunraは、ClaudeやChatGPTなどのアシスタントから画像・動画・音声の各モデルを呼び出せる「Sunra MCP」を発表した。モデル本来のAPI料金と同額で、上乗せ手数料を設けないと説明している。公式投稿にはデモ動画があり、製品ページとMCPサーバーの公開情報が連携方式を紹介している。複数の生成モデルを既存の対話環境から使える点と料金体系が関心を集めた。ただし、料金の一致やツール呼び出しの安定性は本号では実測していない。
\dailyxsource{Sunra (@sunra\_ai)}{https://x.com/sunra_ai/status/2112212015045296395}

\section{米DHS、連邦民間機関にDeepSeek排除を指示}
米国土安全保障省（DHS）が、連邦政府の民間機関に対しDeepSeek製品と、その実装を利用した第三者サービスの削除および新規調達停止を求めたと報じられた。理由には、中国への機微情報の移転やマルウェアの可能性などの安全保障上の懸念が挙げられている。GrokはReutersなどの報道とFederal Registerの資料を根拠として記録した。指示は7月24日付で、10月19日に公表・周知されたという。防衛機関への適用や具体的な実施範囲は、追加確認が必要である。
\dailyxsource{関連報道が共有された投稿 (@AndrewCurran\_)}{https://x.com/AndrewCurran_/status/2112250754358788437}

\section{DeepSeekとHuaweiのAI半導体協業報道}
DeepSeekとHuaweiが、NVIDIAへの依存を減らすためにAIチップの設計やTileLangなどのソフトウェアツールで協力している、との報道がXで再び共有された。投稿はThe Informationの記事を参照し、過去のReuters報道なども背景情報として挙げている。ただし、今回の観測窓内には両社の公式説明や一次資料が十分に示されていない。協業の具体的な範囲や成果は未確定であり、中国のAI技術自立に関する見方も報道・コミュニティの主張として区別する。
\dailyxsource{報道紹介の投稿 (@unusual\_whales)}{https://x.com/unusual_whales/status/2112287169381294235}

\section{Skyfall AI、CAD・BIM向け「SKY-001」を紹介}
Skyfall AIのSam Pasupalak氏は、CAD・BIMなどの設計業務に特化した「Engineering World Model」としてSKY-001を紹介した。投稿では、GPT-6 Astraなどの先端モデルと比べて処理速度が51～125倍、費用が150～210分の1で、IoU指標も高いと主張している。動画デモは共有されたが、Grokの収集では技術報告書や公開リポジトリ、第三者による再現結果を確認できていない。数値は発表者による未検証の比較として扱い、実用性の評価は保留する。
\dailyxsource{Sam Pasupalak (@spisallyouneed)}{https://x.com/spisallyouneed/status/2112242023805964327}

\section{NVIDIA、ドイツでの産業AI連携を紹介}
NVIDIAは、Jensen Huang氏とドイツのMerz首相がベルリンで、ドイツの製造技術とNVIDIAのAIを組み合わせる可能性を議論したと投稿した。対象には製造業や自動車産業が含まれ、現地企業への投資構想にも言及している。産業向けAIの導入やオープンモデルの利用拡大に関連する動きとして注目された。一方、今回のGrok収集だけでは、個別の投資額や契約内容、具体的な実施日程までは確認できない。
\dailyxsource{NVIDIA (@nvidia)}{https://x.com/nvidia/status/2112269224949485860}

\section{エージェント関連の小規模な更新・実演}
主要発表のほか、Reflection AIのBeamを用いて、公開データからニューヨーク市地下鉄の運行状況を示す地図を作成するデモなどが共有された。エージェントが外部データを取り込み、利用者に見える成果物を組み立てる用途への関心がうかがえる。Grokはこの項目について特定の一次X投稿URLを提示していないため、リンクは付けない。観測された小規模なツール更新として扱い、新しい主要モデルの公開とは区別する。
